\documentclass[11pt]{article}
\usepackage[margin=0.8in]{geometry}
\usepackage{amsmath,amssymb,booktabs,microtype,titlesec}
\titlespacing*{\section}{0pt}{6pt}{3pt}
\setlength{\parskip}{3pt}
\title{\vspace{-1.3cm}\textbf{The stylized model}}
\author{}\date{}
\begin{document}\maketitle\vspace{-3.2em}

\section*{Demography and timing}
Overlapping generations. A household lives two periods, young then old, and each
cohort has unit mass. It has one child \emph{while young}; the child leaves at
the end of that period and enters the model as a young adult, so this period's
children are next period's young and the cohort entering is unit mass forever.
The equivalence scale is therefore $\xi$ while young and $1$ when old: the old
carry no child, in the scale or in the objective.

The disorder is realised at the \emph{transition} between the two periods, and
survival over that transition depends on it. It cannot be realised inside the
final period, where death costs nothing because nothing follows it.

\section*{Period 1}
\textbf{State.} Permanent income $y$. An employment draw and, if income-eligible,
a voucher draw give realised income and a rent multiplier $(\text{inc},r)$.
Realised income is floored at $\underline y$. A voucher is means tested at $y_V$,
arrives with probability $\pi_V$, and cuts the tenant's rent to $\kappa_r$ of the
market rent.

\textbf{Choices.} Drug purchase $v\ge0$ at price $p_v$, housing
$h\in\{h_1,h_2\}$ --- the street, which is free and supplies $h_1$, or the one
rental unit, which supplies $h_2$ at rent $p_r h_2$ --- and saving $a'\ge0$,
subject to
\[
  c=\text{inc}-p_v v-\text{rent}(h,r)-a'>0 .
\]

\textbf{Flow utility.} With $\varphi_h=\gamma(1-\theta)$,
$\varphi_c=(1-\gamma)(1-\theta)$, $\varphi_q=\theta$ and
$u(x)=x^{1-\sigma}/(1-\sigma)$,
\[
  X_P=\Big(\tfrac{c}{\xi}\Big)^{\varphi_c}\Big(\tfrac{h}{\xi}\Big)^{\varphi_h}(\bar b+v)^{\varphi_q},
  \qquad
  X_K=\Big(\tfrac{c}{\xi}\Big)^{\varphi_c}\Big(\tfrac{h}{\xi}\Big)^{\varphi_h}\;\bar b^{\,\varphi_q},
\]
\[
  \text{flow}=u(X_P)+\omega_c(\xi-1)\,u(X_K)+\aleph .
\]
$c$ and $h$ are divided by $\xi$ and $v$ is not, because the child consumes none
of the drug: the child's bundle carries the Stone--Geary shifter $\bar b$ alone.
The drug still reaches the child, through the \emph{budget}. It is bought out of
the same resources as consumption and rent, so it lowers the child's $c$ and can
also tip the household out of the rental and onto the street --- and both $c$ and
$h$ are in $X_K$. The child's weight is $\xi-1$, so the household's weights sum
to $\xi$, the same $\xi$ that deflates its bundle. $\aleph$ is added once per
household, not once per member.

\textbf{Onset and mortality.} Onset is a \emph{flat} rate, not a function of how
much is bought: the disorder arrives at the transition with probability $\pi$ if
$v>0$ and with probability zero if $v=0$. Survival is $s_0$ when clean and
$s_1<s_0$ with the disorder. So the drug decision has an extensive margin that
carries all of the risk, and an intensive margin that carries none.

\textbf{The problem.} Write $\mathbf 1_v=\mathbf 1\{v>0\}$, so the onset
probability is $\pi\mathbf 1_v$. Choose $(v,h,a')$ to maximise
\[
  \boxed{\;
  \mathcal V=u(X_P)+\omega_c(\xi-1)u(X_K)+\aleph
  +\beta\Big[(1-\pi\mathbf 1_v)\,s_0W_2(a',y,0)+\pi\mathbf 1_v\,s_1W_2(a',y,1)\Big]
  +\beta\,\omega_c\,\mathbb E\big[V(1,d')\big]\;}
\]
where $V(1,d')$ is the child's own value as a young adult, $d'$ the disorder the
parent ends with, and
\[
  \mathbb E\big[V(1,d')\big]=(1-\pi\mathbf 1_v)V(1,0)+\pi\mathbf 1_v V(1,1)
  \;=\;V(1,0)-\pi\mathbf 1_v\,\Delta(y),
  \qquad \Delta(y)\equiv V(1,0)-V(1,1).
\]
Only the second piece moves with the household's choice, so the drug decision
faces exactly $-\beta\omega_c\pi\mathbf 1_v\Delta(y)$: the child's loss, weighted
by how likely the parent is to acquire the disorder at all.

Survival sits \emph{inside} the expectation over the realised disorder, not
outside it: written as $\phi\,\mathbb{E}W_2$, the disorder becomes something the
household wants wherever $W_2<0$. The $v=0$ branch is solved on its own, with the
clean continuation and the voucher unrevoked, and the household uses if the best
$v>0$ beats it.

\section*{Period 2 --- terminal}
Employment is redrawn, at a worse rate with the disorder, and the voucher is
redrawn subject to the policy rule. Assets are run down and nothing is
bequeathed, so with $\text{cash}=Ra+\text{inc}$ and $\xi=1$,
\[
  W_2(a,y,d)=\max_{h\in\{h_1,h_2\}}\;u\big(c^{\varphi_c}h^{\varphi_h}\bar b^{\,\varphi_q}\big)+\aleph,
  \qquad c=\text{cash}-\text{rent}(h,r).
\]
No drug choice, no child term, no recovery.

\textbf{Policy.} Housing First pays the voucher to every eligible household.
Treatment First withdraws it from any household failing the sobriety condition
\emph{in the period the voucher is paid} --- current use, or a current disorder.
That single difference is the whole experiment.

\section*{The child}
Income is inherited with persistence $\rho$, with the innovation that keeps the
cross-section stationary:
\[
  \log y_c=\rho\log y+(1-\rho)\log\bar y-\delta_d\mathbf 1\{d'=1\}+\varepsilon,
  \qquad \varepsilon\sim N\!\big(0,(1-\rho^2)\sigma_y^2\big).
\]
$\delta_d$ is the only inherited penalty: persistence alone delivers the
homelessness odds ratio, so there is no separate homelessness penalty and the
mixture has two types. The child is then run through the model as an adult at its
own income, which is what gives it a disorder probability at all.

So the two branches of the child's value are
\[
  V(1,0)=\mathbb E\big[V_1(y_c)\,\big|\,\text{no penalty}\big],
  \qquad
  V(1,1)=\mathbb E\big[V_1(y_c)\,\big|\,\text{penalty}\big],
\]
both expectations over the same kernel, and $\Delta(y)=V(1,0)-V(1,1)$. $V_1$
appears on both sides, so this is a fixed point. It is a well behaved one: the
$\aleph$ levels cancel in the difference. There is no survival factor on it ---
the child becomes an adult whether or not the parent lives, so $s_d$ belongs on
$W_2$ and on nothing else.

\section*{Lifetime utility}
$V_1(y)$ is the household's \emph{own} value: the maximised $\mathcal V$ with the
internalised child term netted out, so that a planner counts the child once
rather than once in full and again at $\omega_c$. The
childhood a cohort actually lives, at its parent's realised allocation, is
\[
  W_{\text{child}}=(\xi-1)\big[u(X_K)+\aleph\big],
  \qquad\text{and}\qquad
  W=W_{\text{child}}+\beta V_1 ,
\]
with childhood lived under the policy in force one period earlier and adulthood
under the policy in force now. That asymmetry is what a static comparison sums
away: the cohort the reform lands on has its childhood already behind it.

Consumption equivalents are CRRA, not log. Scaling $c$ by $(1+\lambda)$ in every
state scales both bundles by $(1+\lambda)^{\varphi_c}$, so every $u$ term scales
by $(1+\lambda)^{\varphi_c(1-\sigma)}$ while the $\aleph$ terms do not move:
\[
  \lambda=\Big(1+\tfrac{dV}{U}\Big)^{\frac{1}{\varphi_c(1-\sigma)}}-1,
  \qquad U=V-A,\qquad A=\aleph\Big(1+\beta\big[(1-\pi)s_0+\pi s_1\big]\Big),
\]
which depends on the \emph{level} through $U$, as the log formula never did.

\section*{What identifies what}
\begin{tabular}{@{}ll@{\hskip 3em}ll@{}}
\toprule
$\gamma$ & housing share of non-drug spending & $\theta$ & drug spending share among users\\
$\aleph$ & value of a statistical life & $\pi_{\text{base}}$ & disorder rate\\
$h_1$ or $p_rh_2$ & homelessness rate & $\omega_c$ & the child use gap\\
$\bar b$ & drug participation & $\delta_d$ & the child's SUD odds ratio\\
\bottomrule
\end{tabular}

\medskip
\noindent$\sigma$, $\beta$, $\rho$, $s_0$, the unemployment rates, the unit sizes
and the voucher rule are set rather than calibrated. Participation and the
disorder rate are tied by $\text{sud}=\text{participation}\times\pi_{\text{base}}$,
so they cannot be hit independently. $\omega_c$ acts in two places, the flow and
the continuation, and the continuation half is $\Delta(y)$.
\end{document}
